\section{Introduction}

Le projet de lecture dirigée en mathématiques et statistique offre aux étudiantes et étudiants de premier cycle l'occasion d'explorer, en profondeur, un sujet avancé qui n'est habituellement pas abordé dans le cadre du baccalauréat. Cet accompagnement individualisé se déroule sous la supervision d'un mentor ou d'une mentore inscrit(e) aux cycles supérieurs, dans un cadre qui favorise l'autonomie, la rigueur et la curiosité mathématique. Il se veut amusant et amical dans le but de faciliter les interactions entre les différents membres.\\
De plus, ce projet représente également une occasion pour les étudiant(e)s aux cycles supérieurs d'enrichir leur expérience de l'enseignement et de l'encadrement, tout en contribuant activement à la vie scientifique du département. Il vise également à tisser des liens entre les différentes cohortes étudiantes et à stimuler un environnement mathématique dynamique. \\ 
Pour atteindre ces objectifs, le projet s'appuie sur une équipe organisatrice ayant déjà fait ses preuves dans la coordination d'activités scientifiques similaires. \\
En effet, Emanuele Ronda fût organisateur dans le cadre des séminaires CIRGET Junior ainsi que de l'école d'été ISM \'Ecole Découverte : Géométrie de K\"{a}hler, géométrie algébrique et arithmétique, hyperbolicité ; cela lui a permis de gagner de l'expérience dans la logistique et la gestion d'un budget. Pascal Guiffo Kamwa sera également organisateur des séminaires CIRGET Junior, lui permettant d'obtenir les mêmes types d'expériences. Finalement, Ismael El Yassini a déjà été mentor pour le projet DREAMS dans son édition passée, donc il rajoutera la connaissance spécifique liée à ce projet. Cette expérience combinée couvre l'ensemble des aspects du projet — logistique, encadrement pédagogique et continuité avec l'édition précédente — ce qui réduit les risques associés à une première mise en œuvre.
